\chapter{Skin and Soft Tissue}

\section*{Patient focus}
This chapter explains the examination or clinical problem in accessible language. Patients should use it to prepare questions and understand the purpose, experience, benefits, and limitations of imaging; local instructions and the official report take priority.


\label{ch:ch55_skin_soft_tissue}

\section{What does soft tissue imaging cover?}
Soft tissue imaging evaluates structures outside the skeleton---muscles, subcutaneous fat, fascia, skin, and neurovascular bundles. It is used for characterising lumps, assessing infection, staging tumours, and evaluating trauma.

\section{When is imaging ordered?}
\begin{itemize}
    \item \textbf{Palpable lump:} Is it a lipoma, cyst, abscess, or something more sinister?
    \item \textbf{Suspected infection:} Cellulitis, abscess, or necrotising fasciitis.
    \item \textbf{Soft tissue swelling without clear cause:} Assess extent and underlying cause.
    \item \textbf{Known or suspected sarcoma:} Characterisation and staging.
    \item \textbf{Trauma:} Muscle tear, haematoma, or compartment syndrome.
    \item \textbf{Foreign body:} Retained glass, wood, or metal.
\end{itemize}

\section{Modality selection}

\begin{table}[h]
\centering
\caption{Imaging modalities for soft tissue}
\label{tab:st_modality}
\begin{tabular}{lll}
\hline
\textbf{Modality} & \textbf{Best for} & \textbf{Limitations} \\
\hline
CT & Gas in tissues (necrotising fasciitis), calcification, bone involvement & Radiation; less soft tissue contrast than MRI \\
MRI & Extent of tumours, fascial planes, neurovascular, marrow & Cost, time, contraindications \\
Ultrasound & Superficial masses, fluid, vascularity, foreign bodies & Operator dependent; deep lesions missed \\
X-ray & Soft tissue calcification, gas, foreign bodies (metal/glass) & Low sensitivity for most soft tissue pathology \\
\hline
\end{tabular}
\end{table}

\section{Common soft tissue masses}

\begin{table}[h]
\centering
\caption{Common soft tissue masses}
\label{tab:st_masses}
\begin{tabular}{llll}
\hline
\textbf{Mass} & \textbf{US features} & \textbf{MRI features} & \textbf{Key point} \\
\hline
Abscess & Complex fluid, rim enhancement & Ring-enhancing collection & Clinical signs of infection \\
Epidermoid cyst & Hypoechoic, may have tract & T1 bright (proteinaceous) & Central punctum on skin \\
Ganglion & Anechoic, well-defined & T2 bright, no solid component & Often near joints/tendons \\
Lipoma & Echogenic, compressible & Follows fat on all sequences & Most common soft tissue mass \\
Nodular fasciitis & Rapid growth, solid & Intermediate signal & Can mimic sarcoma; benign \\
Sarcoma & Solid, vascular, deep & Heterogeneous, enhancing & Urgent referral if suspected \\
\hline
\end{tabular}
\end{table}

\begin{pearlbox}[Lipoma vs liposarcoma]
A subcutaneous lesion following fat signal on all MRI sequences with thin, regular septa is almost certainly a lipoma. Suspect liposarcoma if: the lesion is deep to fascia, $>$10~cm, has thick ($>$2~mm) or nodular septa, has non-fatty components, or is growing. These features warrant urgent sarcoma referral.
\end{pearlbox}

\section{Skin thickening and subcutaneous oedema}
Skin thickening (measured $>$2~mm on ultrasound, or visible thickening on MRI/CT) can indicate:
\begin{itemize}
    \item \textbf{Cellulitis:} Diffuse skin and subcutaneous thickening with oedema and stranding.
    \item \textbf{Lymphoedema:} Non-pitting oedema with skin thickening, often with ``honeycomb'' subcutaneous pattern on MRI.
    \item \textbf{Venous insufficiency:} Skin thickening and subcutaneous oedema, typically in the lower limbs, with varicose veins.
    \item \textbf{Inflammatory dermatoses:} Eczema, psoriasis, dermatomyositis.
    \item \textbf{Malignancy:} Inflammatory breast cancer (peau d'orange), lymphangitis carcinomatosa.
\end{itemize}

\section{Cellulitis and abscess}
\begin{itemize}
    \item \textbf{Cellulitis on ultrasound:} Thickened, hypoechoic skin and subcutaneous tissue with reticular pattern. Increased Doppler flow. No defined fluid collection.
    \item \textbf{Abscess on ultrasound:} Well-defined, hypoechoic or anechoic fluid collection with internal debris, posterior enhancement, and peripheral hyperaemia. Size and depth guide management (incision and drainage vs aspiration).
    \item \textbf{Cellulitis on MRI:} Skin thickening, subcutaneous fat stranding (high T2/STIR signal), and enhancement. A rim-enhancing fluid collection indicates abscess.
\end{itemize}

\section{Necrotising fasciitis: MRI findings}
Necrotising fasciitis is a life-threatening surgical emergency. MRI is the most sensitive imaging modality:
\begin{itemize}
    \item \textbf{T2/STIR:} High signal along fascial planes (fascial oedema).
    \item \textbf{T1 post-contrast:} Thickening and enhancement of the deep fascia. Non-enhancing areas indicate necrosis.
    \item \textbf{Gas in soft tissues:} Seen as signal voids on MRI, or as lucencies on CT. Gas may not always be present early.
    \item \textbf{Key finding:} Fluid between the muscle and fascia, or between fascial layers.
\end{itemize}

\begin{pearlbox}[Clinical suspicion trumps imaging]
Necrotising fasciitis is a clinical diagnosis. Imaging should not delay surgical debridement. If there is clinical concern (pain out of proportion, crepitus, systemic toxicity, rapidly spreading erythema), proceed to the operating theatre---imaging is adjunctive, not diagnostic.
\end{pearlbox}

\section{Compartment syndrome}
Compartment syndrome occurs when tissue pressure within a fascial compartment exceeds perfusion pressure, leading to ischaemia.
\begin{itemize}
    \item \textbf{Clinical:} The 6 P's---pain (out of proportion), pallor, pulselessness, paraesthesia, paralysis, and poikilothermia.
    \item \textbf{Imaging role:} MRI shows compartment oedema (high T2 signal in muscles), increased intermuscular pressure, and venous congestion. CT angiography may show reduced arterial flow.
    \item \textbf{Important:} Compartment syndrome is primarily a clinical diagnosis. Intracompartmental pressure measurement (needle manometry) confirms the diagnosis. Imaging is rarely required and must not delay fasciotomy.
\end{itemize}

\section{MRI protocol for soft tissue masses}
A dedicated soft tissue MRI protocol includes:
\begin{enumerate}
    \item \textbf{T1-weighted:} Anatomy, fat identification, marrow assessment.
    \item \textbf{T2-weighted fat-saturated or STIR:} Oedema, fluid, tumour characterisation.
    \item \textbf{T1 post-gadolinium fat-saturated:} Enhancement pattern (solid vs cystic, rim vs homogeneous).
    \item \textbf{DWI:} High-grade tumours show restricted diffusion.
    \item \textbf{Field of view:} Must include the entire lesion plus adjacent joint above and below.
\end{enumerate}

\section{When to suspect sarcoma: referral criteria}
Refer urgently if a soft tissue mass is:
\begin{itemize}
    \item $>$5~cm in maximum diameter
    \item Deep to the fascia
    \item Increasing in size
    \item Associated with pain (especially night pain)
    \item Solid and vascular on Doppler ultrasound
    \item Has imaging features of malignancy (heterogeneous enhancement, restricted diffusion, irregular margins)
\end{itemize}

\section{Reporting checklist}
A structured soft tissue report should include:
\begin{enumerate}
    \item Location (anatomical compartment, superficial vs deep)
    \item Size (3 dimensions)
    \item Signal/echogenicity characteristics
    \item Enhancement pattern
    \item Relationship to fascia, muscle, bone, and neurovascular structures
    \item Vascularity (Doppler)
    \item Evidence of infection (gas, abscess, fascial involvement)
    \item Interval change (if prior imaging available)
    \item Impression with clear management recommendation
\end{enumerate}

\clearpage
